\documentclass[10pt,a4paper]{article}
\usepackage{bold}

\title{$\sigma_0$ \,:\, large scale quant foundation models}
\author{Aaron Rose}
\started{5 September 2026}

\begin{document}
\maketitle
\boldcontents

% ================================================================
\section{Data}

\subsection{The book}

A snapshot of the order book at time $t$ is
%
\begin{equation}
S_t \;=\;
\begin{bmatrix}
a_1 & u_1 & b_1 & w_1 \\[1pt]
\vdots & \vdots & \vdots & \vdots \\[1pt]
a_{10} & u_{10} & b_{10} & w_{10}
\end{bmatrix}
\;\in\; \mathbb{R}^{10\times 4},
\label{eq:snapshot}
\end{equation}
%
with $a_j$ and $b_j$ the ask and bid price at level $j$, and $u_j$ and $w_j$
the depth at that price. Prices are integers in units of $10^{-4}$ dollars, so
one cent is $100$. Figure~\ref{fig:book} shows one snapshot.

% ---------------------------------------------------------------- fig 1
\begin{figure}[H]
\centering
\begin{tikzpicture}[x=1pt,y=1pt,font=\footnotesize,color=Ink]

\node[anchor=east] at (-6,88) {price};
\node[anchor=east] at (104,88) {size};

% ---- asks, best at the bottom, price increasing upward
\foreach \i/\p/\s in {1/1010100/200, 2/1010200/1200, 3/1010300/400,
                      4/1010400/900, 5/1010500/1500, 6/1010600/700,
                      7/1010700/1100, 8/1010800/300, 9/1010900/1800,
                      10/1011000/600} {
  \pgfmathsetmacro{\y}{6+(\i-1)*7.8}
  \pgfmathsetmacro{\w}{\s/22}
  \fill[Noisy!45] (0,\y) rectangle (\w,\y+4.8);
  \node[anchor=east] at (-6,\y+2.4) {\p};
  \node[anchor=east] at (104,\y+2.4) {\s};
}

% ---- bids, best at the top, price decreasing downward
\foreach \i/\p/\s in {1/1009900/600, 2/1009800/1000, 3/1009700/150,
                      4/1009600/1300, 5/1009500/800, 6/1009400/400,
                      7/1009300/1600, 8/1009200/900, 9/1009100/500,
                      10/1009000/1200} {
  \pgfmathsetmacro{\y}{-12-(\i-1)*7.8}
  \pgfmathsetmacro{\w}{\s/22}
  \fill[Data!45] (0,\y) rectangle (\w,\y+4.8);
  \node[anchor=east] at (-6,\y+2.4) {\p};
  \node[anchor=east] at (104,\y+2.4) {\s};
}

% ---- the untraded band between the two best quotes
\draw[Ink!30,line width=0.3pt] (0,6) -- (84,6);
\draw[Ink!30,line width=0.3pt] (0,-6) -- (84,-6);
\draw[<->,Acc,line width=0.5pt] (14,6) -- (14,-6);
\node[anchor=west,color=Acc,fill=white,inner sep=1.5pt] at (18,0) {$a_1-b_1=200$};

% ---- braces naming the two sides
\draw[decorate,decoration={brace,amplitude=4pt},Noisy,line width=0.5pt]
      (110,81) -- (110,6);
\node[anchor=west,color=Noisy] at (117,43.5) {asks $a_1 \ldots a_{10}$};
\draw[decorate,decoration={brace,amplitude=4pt},Data,line width=0.5pt]
      (110,-6) -- (110,-82);
\node[anchor=west,color=Data] at (117,-44) {bids $b_1 \ldots b_{10}$};

\end{tikzpicture}
\caption{One snapshot, ten levels a side.}
\label{fig:book}
\end{figure}

A day is $N$ such snapshots,
%
\begin{equation}
W \;=\;
\bigl[\;
\underbrace{\delta \;\;\; t_{\mathrm{s}} \;\;\; t_{\mathrm{n}}}_{\text{metadata},\ 3}
\;\;{\color{Ink!45}\rule[-0.55ex]{0.4pt}{2.1ex}}\;\;
\underbrace{a_1 \;\; u_1 \;\; b_1 \;\; w_1 \;\;\cdots\;\;
            a_{10} \;\; u_{10} \;\; b_{10} \;\; w_{10}}_{\text{snapshot},\ 40}
\;\bigr]
\;\in\; \mathbb{R}^{N\times 43},
\label{eq:wide}
\end{equation}
%
where $\delta$ is the change in the mid, and $t_{\mathrm{s}}$ and
$t_{\mathrm{n}}$ are the seconds and nanoseconds of the event's timestamp, kept
apart so that neither loses integer precision. Converting the seconds to
nanoseconds and adding the remainder gives the time of the event,
%
\begin{equation}
\tau \;=\; 10^{9}\,t_{\mathrm{s}} + t_{\mathrm{n}} ,
\label{eq:tau}
\end{equation}
%
which is what places it on the grid in \eqref{eq:ffill}.

\subsection{Discrete time}

The book has a state at nanosecond resolution. It is sampled at $1$ ms
intervals,
%
\begin{equation}
t_k \;=\; t_0 + k\Delta, \qquad k = 0,1,2,\ldots, \qquad \Delta = 1\ \text{ms},
\label{eq:grid}
\end{equation}
%
with $F_k$ the book at $t_k$,
%
\begin{equation}
F_k \;=\; S_{t_k} \;\in\; \mathbb{R}^{10\times 4}.
\label{eq:ffill}
\end{equation}
%
% ---------------------------------------------------------------- fig clock
\begin{figure}[H]
\centering
\begin{tikzpicture}[x=1pt,y=1pt,font=\footnotesize,color=Ink]

% ---- sample times
\foreach \x/\k in {20/0, 60/1, 100/2, 140/3, 180/4, 220/5, 260/6} {
  \draw[Ink!35,line width=0.3pt] (\x,-14) -- (\x,32);
  \node[anchor=south] at (\x,33) {$F_{\k}$};
  \node[anchor=north] at (\x,-16) {$t_{\k}$};
}

% ---- the tape, buys above the line and sells below
\draw[Ink,line width=0.5pt] (-6,0) -- (278,0);
\node[anchor=west] at (281,0) {time};
\foreach \x in {14,35,52,72,112,130,168,175,232,255}
  \draw[Data,line width=1.2pt] (\x,0) -- (\x,6);
\foreach \x in {8,48,55,95,118,155,172,250}
  \draw[Noisy,line width=1.2pt] (\x,0) -- (\x,-6);
\draw[Data,line width=1.2pt] (285,-15) -- (285,-9);
\node[anchor=west] at (288,-12) {buy};
\draw[Noisy,line width=1.2pt] (285,-27) -- (285,-21);
\node[anchor=west] at (288,-24) {sell};

% ---- which event each frame takes its state from
\foreach \e/\g in {14/20, 55/60, 95/100, 130/140, 175/180, 255/260} {
  \draw[Acc,->,line width=0.5pt] (\e,8) |- (\g,14);
}
\draw[Acc,->,dashed,line width=0.5pt] (181,22) -- (219,22);
\node[anchor=south,color=Acc] at (200,23) {repeat};

\end{tikzpicture}
\caption{Each frame holds the book as the last event left it.}
\label{fig:clock}
\end{figure}

This is done over the ten level book data, one trading day of one ticker at a
time. The corpus covers the S\&P 500 from 2022 to 2025, and the first state
model runs used GOOG.

The frames are cut into consecutive windows of $L=512$,
%
\begin{equation}
F \;=\;
\begin{bmatrix} F_0 \\[1pt] \vdots \\[1pt] F_{L-1} \end{bmatrix}
\;\in\; \mathbb{R}^{512\times 40},
\label{eq:window}
\end{equation}
%
and the day is these windows end to end,
%
\begin{equation}
\underbrace{F^{(1)},\;\; F^{(2)},\;\; F^{(3)},\;\; \ldots,\;\; F^{(M)}}
           _{\text{one day},\ M \approx 45{,}700} .
\label{eq:stream}
\end{equation}
%
Given the first $n$ rows of a window, the remaining $L-n$ are generated in one
pass. Training draws $n$ from $\{0,\ldots,256\}$, so between $256$ and $512$ ms
is generated.

% ---------------------------------------------------------------- fig stream
\begin{figure}[H]
\centering
\begin{tikzpicture}[x=1pt,y=1pt,font=\footnotesize,color=Ink]

% ---- legend
% bounding box forced symmetric about the bars, so they sit centred on the page
\path (-43,0) (339,0);
\fill[Data!75] (93.5,55) rectangle (102.5,64);
\node[anchor=west] at (105.5,59.5) {given};
\fill[Data!12] (145.5,54) rectangle (155.5,65);
\draw[pattern=crosshatch dots,pattern color=Data,draw=Ink!50,line width=0.3pt]
      (145.5,54) rectangle (155.5,65);
\node[anchor=west] at (158.5,59.5) {generated};

% ---- one window, three amounts of context
\foreach \n/\g/\y/\out in {0/48/34/512, 128/98/17/384, 256/148/0/256} {
  \node[anchor=east] at (42,\y+7) {$n=\n$};
  \fill[Data!75] (48,\y) rectangle (\g,\y+14);
  \fill[Data!12] (\g,\y) rectangle (248,\y+14);
  \draw[pattern=crosshatch dots,pattern color=Data,draw=none] (\g,\y) rectangle (248,\y+14);
  \draw[Ink!50,line width=0.3pt] (48,\y) rectangle (248,\y+14);
  \node[anchor=west] at (254,\y+7) {generate $\out$ ms};
}

\end{tikzpicture}
\caption{Whatever is given, the rest of the window is generated.}
\label{fig:stream}
\end{figure}

\subsection{The chart transform}

The book states are encoded into a different but lossless representation. The
model reads and writes these coordinates, and the book is recovered by decoding
at the end. With the mid
%
\begin{equation}
m_k \;=\; \tfrac{1}{2}\bigl(a_1+b_1\bigr)
\label{eq:mid}
\end{equation}
%
at frame $k$, Table~\ref{tab:chart} gives every entry.

\begin{table}[H]
\centering\small
\begin{tabular}{cll}
\toprule
 & in words & expression \\
\midrule
$s$   & the log spread                    & $\log(a_1-b_1)$ \\
$d$   & change in the mid over one frame  & $m_k - m_{k-1}$, zero at $k=0$ \\
$g_j$ & log distance between ask levels $j$ and $j+1$ & $\log(a_{j+1}-a_j)$ \\
$h_j$ & log distance between bid levels $j$ and $j+1$ & $\log(b_j-b_{j+1})$ \\
$u_j$ & depth at ask level $j$            & unchanged \\
$w_j$ & depth at bid level $j$            & unchanged \\
\bottomrule
\end{tabular}
\caption{The entries of a chart frame.}
\label{tab:chart}
\end{table}

Distances run $j=1,\ldots,9$ and depths $j=1,\ldots,10$, and each logarithm
takes a floor of $1$, the price unit. A snapshot and a window then read
%
\begin{equation}
C_k \;=\;
\begin{bmatrix}
s & u_1 & d & w_1 \\[1pt]
g_1 & u_2 & h_1 & w_2 \\[1pt]
\vdots & \vdots & \vdots & \vdots \\[1pt]
g_9 & u_{10} & h_9 & w_{10}
\end{bmatrix}
\in \mathbb{R}^{10\times 4},
\qquad
C \;=\;
\begin{bmatrix} C_0 \\[1pt] \vdots \\[1pt] C_{L-1} \end{bmatrix}
\in \mathbb{R}^{512\times 40},
\label{eq:chart}
\end{equation}
%
the same shapes as $S_t$ and $F$. Figure~\ref{fig:chart} marks the three
distances on the book.

% ---------------------------------------------------------------- fig chart
\begin{figure}[H]
\centering
\begin{tikzpicture}[x=1pt,y=1pt,font=\footnotesize,color=Ink]

% ---- three levels a side, the bar centre carrying the price
\foreach \i/\s in {1/200, 2/1200, 3/400} {
  \pgfmathsetmacro{\y}{10+(\i-1)*13}
  \pgfmathsetmacro{\w}{\s/26}
  \fill[Noisy!45] (0,\y-2.5) rectangle (\w,\y+2.5);
  \node[anchor=east,color=Noisy] at (-4,\y) {$a_\i$};
}
\foreach \i/\s in {1/600, 2/1000, 3/150} {
  \pgfmathsetmacro{\y}{-10-(\i-1)*13}
  \pgfmathsetmacro{\w}{\s/26}
  \fill[Data!45] (0,\y-2.5) rectangle (\w,\y+2.5);
  \node[anchor=east,color=Data] at (-4,\y) {$b_\i$};
}

% ---- the mid
\draw[Ink!35,densely dotted] (0,0) -- (60,0);
\node[anchor=east] at (-4,0) {$m_k$};

% ---- the three marked distances
\draw[<->,Acc,line width=0.5pt] (70,10) -- (70,-10);
\node[anchor=west,color=Acc] at (74,0) {$a_1-b_1$};
\draw[<->,Acc,line width=0.5pt] (120,23) -- (120,10);
\node[anchor=west,color=Acc] at (124,16.5) {$a_2-a_1$};
\draw[<->,Acc,line width=0.5pt] (120,-10) -- (120,-23);
\node[anchor=west,color=Acc] at (124,-16.5) {$b_1-b_2$};

% ---- the resulting entries
\draw[->,Ink,line width=0.5pt] (176,0) -- (196,0);
\node[anchor=south] at (186,2) {encode};
\node[anchor=west] at (206,22) {$s\;=\;\log(a_1-b_1)$};
\node[anchor=west] at (206,9)  {$d\;=\;m_k-m_{k-1}$};
\node[anchor=west] at (206,-6) {$g_1\;=\;\log(a_2-a_1)$};
\node[anchor=west] at (206,-20){$h_1\;=\;\log(b_1-b_2)$};

\end{tikzpicture}
\caption{Each price becomes a log distance.}
\label{fig:chart}
\end{figure}

Decoding inverts this,
%
\begin{equation}
m_k \;=\; m_0 + \sum_{i\le k} d_i, \qquad
a_1 \;=\; m_k + \tfrac{1}{2}e^{s}, \qquad
b_1 \;=\; m_k - \tfrac{1}{2}e^{s},
\label{eq:decodetouch}
\end{equation}
%
\begin{equation}
a_{j+1} \;=\; a_j + e^{g_j}, \qquad
b_{j+1} \;=\; b_j - e^{h_j} .
\label{eq:decodeladder}
\end{equation}
%
with the anchor $m_0$, the mid of the first frame, the only absolute price
kept. In single precision the round trip returns integer prices exactly.

Why the transform is used.

\begin{itemize}\setlength{\itemsep}{1pt}
\item Impossible books become impossible. A crossed book, with buys above
      sells, has no representation in these coordinates. BookDiT reports
      $54.4$ percent of frames crossed when the same model is trained in raw
      coordinates, against $0.000$ in this one \cite{bookdit}.
\item Fewer digits, and a stationary target. Raw prices are of order $10^6$ and
      drift over months, where their log distances are of order $10$ and do
      not.
\end{itemize}

\subsection{The input tensor}

Each row of a window is a frame, the book at one instant, and each column is a
channel, a single quantity followed across the whole window.
Figure~\ref{fig:grid} shows the two readings.

% ---------------------------------------------------------------- fig grid
\begin{figure}[H]
\centering
\begin{tikzpicture}[x=1pt,y=1pt,font=\footnotesize,color=Ink]

% ---- the shaded column and row
\fill[Noisy!25] (45,0) rectangle (60,100);
\fill[Data!25]  (0,55) rectangle (150,65);
\fill[Acc!30]   (45,55) rectangle (60,65);

% ---- the grid
\foreach \x in {0,15,...,150} \draw[Ink!30,line width=0.3pt] (\x,0) -- (\x,100);
\foreach \y in {0,10,...,100} \draw[Ink!30,line width=0.3pt] (0,\y) -- (150,\y);
\draw[Ink!60,line width=0.5pt] (0,0) rectangle (150,100);

% ---- what each direction is
\draw[<->,Ink!70,line width=0.4pt] (0,110) -- (150,110);
\node[anchor=south] at (75,112) {$40$ channels};
\draw[<->,Ink!70,line width=0.4pt] (-10,0) -- (-10,100);
\node[anchor=east] at (-14,50) {$512$ frames};

\draw[Noisy,->,line width=0.5pt] (52.5,-16) -- (52.5,-3);
\node[anchor=north,color=Noisy] at (52.5,-18) {a channel};
\draw[Data,->,line width=0.5pt] (166,60) -- (153,60);
\node[anchor=west,color=Data] at (169,60) {a frame};

\end{tikzpicture}
\caption{A window read down its columns and across its rows.}
\label{fig:grid}
\end{figure}

The forty channels differ by orders of magnitude, a log spread near five
sitting beside a depth in the thousands, so each is standardised separately
over the training days,
%
\begin{equation}
\widetilde{C} \;=\; \frac{C-\mu}{\sigma}, \qquad \mu,\ \sigma \;\in\; \mathbb{R}^{40},
\label{eq:norm}
\end{equation}
%
with one mean and one standard deviation per channel. Both are stored inside
every checkpoint, so that generation reproduces the scale the model was trained
on. Windows are then stacked into batches, and the tensor the network is given
is
%
\begin{equation}
X \;=\;
\begin{bmatrix} \widetilde{C}^{(1)} \\[1pt] \vdots \\[1pt] \widetilde{C}^{(B)} \end{bmatrix}
\;\in\; \mathbb{R}^{32\times 512\times 40}.
\label{eq:tensor}
\end{equation}
%
The logarithm of one plus the count of events in each interval can be appended
as a $41$st channel, off by default.

\begin{table}[H]
\centering\small
\begin{tabular}{llr}
\toprule
quantity & symbol & default \\
\midrule
frame interval   & $\Delta$ & $1$ ms \\
levels           &          & $10$ \\
channels         &          & $40$ \\
window length    & $L$      & $512$ frames, $512$ ms \\
maximum context  & $n$      & $256$ frames \\
batch            & $B$      & $32$ \\
\bottomrule
\end{tabular}
\caption{Defaults.}
\label{tab:defaults}
\end{table}

One training example is half a second of book on an axis of time and an axis
of channel, the same object a video model denoises.


\clearpage
% ================================================================
\section{Pretraining}

\subsection{The objective}

Training consumes the tensor of Section~1.4. One row holds one frame, and its
forty entries are the quantities of Table~\ref{tab:chart} ordered by book level,
%
\begin{equation}
\bigl(\;
\underbrace{s \;\; u_1 \;\; d \;\; w_1}_{\text{level }1}
\;\;
\underbrace{g_1 \;\; u_2 \;\; h_1 \;\; w_2}_{\text{level }2}
\;\;\cdots\;\;
\underbrace{g_9 \;\; u_{10} \;\; h_9 \;\; w_{10}}_{\text{level }10}
\;\bigr) \;\in\; \mathbb{R}^{40},
\label{eq:row}
\end{equation}
%
where each entry is standardised by \eqref{eq:norm}. A window is $512$
consecutive rows and a batch is $32$ windows,
%
\begin{equation}
X \;\in\; \mathbb{R}^{32\times 512\times 40},
\label{eq:data}
\end{equation}
%
each batch supplying one gradient step. Figure~\ref{fig:tensor} shows a batch as
a stack of windows.

% ---------------------------------------------------------------- fig tensor
\begin{figure}[H]
\centering
\begin{tikzpicture}[x=1pt,y=1pt,font=\footnotesize,color=Ink]

\foreach \i in {2,1} {
  \pgfmathsetmacro{\ox}{\i*22}
  \pgfmathsetmacro{\oy}{\i*16}
  \fill[white] (\ox,\oy) rectangle (\ox+80,\oy+100);
  \draw[Ink!55,line width=0.6pt]
    (\ox+6,\oy) -- (\ox,\oy) -- (\ox,\oy+100) -- (\ox+6,\oy+100);
  \draw[Ink!55,line width=0.6pt]
    (\ox+74,\oy) -- (\ox+80,\oy) -- (\ox+80,\oy+100) -- (\ox+74,\oy+100);
}
\fill[Wash!10] (0,0) rectangle (80,100);
\foreach \x in {8,16,...,72} \draw[Ink!20,line width=0.3pt] (\x,0) -- (\x,100);
\foreach \y in {10,20,...,90} \draw[Ink!20,line width=0.3pt] (0,\y) -- (80,\y);
\draw[Ink!80,line width=0.8pt] (6,0) -- (0,0) -- (0,100) -- (6,100);
\draw[Ink!80,line width=0.8pt] (74,0) -- (80,0) -- (80,100) -- (74,100);

\node[anchor=south west] at (0,102) {$\widetilde{C}^{(1)}$};
\node[anchor=south west] at (22,118) {$\cdots$};
\node[anchor=south west] at (44,134) {$\widetilde{C}^{(32)}$};

\draw[<->,Ink!70,line width=0.4pt] (0,-9) -- (80,-9);
\node[anchor=north] at (40,-12) {$40$ channels};
\draw[<->,Ink!70,line width=0.4pt] (-10,0) -- (-10,100);
\node[anchor=east] at (-13,50) {$512$ frames};

\end{tikzpicture}
\caption{The batch tensor $X$, thirty two windows of $512$ frames by $40$ channels.}
\label{fig:tensor}
\end{figure}

Each batch comprises thirty two independent draws,
%
\begin{equation}
X^{(b)} \sim p, \quad
\epsilon^{(b)} \sim \mathcal{N}(0,I), \quad
t_b \sim \mathcal{U}[0,1], \quad
n_b \sim \mathcal{U}\{0,\ldots,256\},
\label{eq:draws}
\end{equation}
%
where $p$ is the distribution of the windows of Section~1 and the entries of
$\epsilon^{(b)}$ are independent. The window and its noise
are linearly interpolated at the level drawn for that slot,
%
\begin{equation}
x^{(b)} \;=\; (1-t_b)\,X^{(b)} \;+\; t_b\,\epsilon^{(b)},
\label{eq:interp}
\end{equation}
%
which equals $X^{(b)}$ at $t=0$ and $\epsilon^{(b)}$ at $t=1$.

\begin{table}[H]
\centering\small
\begin{tabular}{lrrrrl}
\toprule
 & $\tilde{s}$ & $\tilde{u}_1$ & $\tilde{d}$ & $\tilde{w}_1$ & \\
\midrule
$X$        &  0.42 & $-1.13$ &  0.05 &  0.88 & $t=0$ \\
           &  0.24 & $-0.41$ & $-0.12$ &  0.69 & $t=0.25$ \\
           &  0.06 &  0.32 & $-0.29$ &  0.51 & $t=0.5$ \\
           & $-0.13$ &  1.05 & $-0.45$ &  0.32 & $t=0.75$ \\
$\epsilon$ & $-0.31$ &  1.77 & $-0.62$ &  0.14 & $t=1$ \\
\midrule
$\epsilon-X$ & $-0.73$ & 2.90 & $-0.67$ & $-0.74$ & the target \\
\bottomrule
\end{tabular}
\caption{Four of the forty standardised channels of one frame, its noise, and
their interpolants.}
\label{tab:mixing}
\end{table}

The interpolant traverses the segment from $X^{(b)}$ to $\epsilon^{(b)}$ in unit
time, and its velocity
%
\begin{equation}
\frac{\mathrm{d}x^{(b)}}{\mathrm{d}t} \;=\; \epsilon^{(b)} - X^{(b)}
\label{eq:velocity}
\end{equation}
%
is the regression target.

The network takes $x^{(b)}$ and $t_b$ as arguments and predicts the velocity.
The minimiser of the squared error is the conditional mean,
%
\begin{equation}
v^{*}(x,t) \;=\;
\operatorname*{argmin}_{v}\;
\mathbb{E}\bigl\|(\epsilon-X) - v(x_t,t)\bigr\|^2
\;=\; \mathbb{E}\bigl[\,\epsilon-X \;\bigm|\; x_t = x\,\bigr],
\label{eq:expectation}
\end{equation}
%
the average velocity of the interpolants passing through $x$ at time $t$
\cite{rectified}. Integrating this field carries noise to samples distributed as
$p$, and Section~3 gives the solver.

% ---------------------------------------------------------------- fig vector
\begin{figure}[H]
\centering
\begin{tikzpicture}[x=1pt,y=1pt,font=\footnotesize,color=Ink]

\node[anchor=west,color=Ink!60] at (16,72) {$\mathbb{R}^{512\times 40}$};

% ---- the segment between a window and its noise
\draw[Ink!35,dashed,line width=0.5pt] (20,25) -- (170,70);
\fill[Data]  (20,25) circle (2.2pt);
\node[anchor=east,color=Data] at (16,25) {$X$};
\fill[Noisy] (170,70) circle (2.2pt);
\node[anchor=west,color=Noisy] at (174,70) {$\epsilon$};

% ---- the same displacement, wherever you stand on it
\draw[Ink!45,->,line width=0.6pt] (42.5,31.75) -- (67.2,39.2);
\draw[Ink!45,->,line width=0.6pt] (128,57.4) -- (152.7,64.8);
\draw[Acc,->,line width=1.2pt] (87.5,45.25) -- (112.2,52.7);
\node[anchor=north west,color=Acc] at (110,51) {$\epsilon-X$};
\fill[Ink] (87.5,45.25) circle (1.8pt);
\node[anchor=south east,inner sep=1.5pt] at (86,46) {$x_t$};

\end{tikzpicture}
\caption{The segment from $X$ to $\epsilon$ in $\mathbb{R}^{512\times 40}$, and its constant velocity.}
\label{fig:vector}
\end{figure}

The field is modelled by a network $v_\theta$ with parameters
$\theta\in\mathbb{R}^{P}$. Writing $\|\cdot\|^2$ for the mean squared entry of
its argument, the batch objective is
%
\begin{equation}
L(\theta) \;=\; \frac{1}{32}\sum_{b=1}^{32}
\bigl\|\,\bigl(\epsilon^{(b)}-X^{(b)}\bigr)
- v_\theta\bigl(x^{(b)},t_b\bigr)\bigr\|^2 .
\label{eq:loss}
\end{equation}
%
The draws are held fixed, so $L$ depends on $\theta$ alone, and it is minimised
by descent. The output layer of $v_\theta$ is initialised at zero, so
$v_\theta$ vanishes identically before training. Section~2.2 conditions the
model on part of each window, which replaces \eqref{eq:loss} with the objective
actually used.


\subsection{Conditioning}

Sampling from $p$ produces $512$ ms of a book resembling recorded ones but
bearing no relation to any particular book.
A simulator starts from the book as it stands and asks what follows, so what is
required is the distribution of continuations of a given beginning. Each window
is therefore drawn with a context length
%
\begin{equation}
n_b \;\sim\; \mathcal{U}\{0,\ldots,256\},
\label{eq:ctx}
\end{equation}
%
fixing how much of it is given to the network and how much it must produce, and
the model samples from
%
\begin{equation}
p\bigl(X_{\ge n_b} \mid X_{< n_b}\bigr),
\label{eq:conddist}
\end{equation}
%
where $X_{<n}$ and $X_{\ge n}$ are the first $n$ frames of a window and the
remainder.

The given frames are supplied on their own channels beside the interpolant,
together with an indicator of which frames they are, giving $81$ input channels
against the $40$ of the unconditional case. Writing $n$ for $n_b$, the input is
%
\begin{equation}
\left[
\begin{array}{ccc}
x_0 & X_0 & 0 \\[1pt]
\vdots & \vdots & \vdots \\[1pt]
x_{n-1} & X_{n-1} & 0 \\[1pt]
x_n & 0 & 1 \\[1pt]
\vdots & \vdots & \vdots \\[1pt]
x_{511} & 0 & 1
\end{array}
\right]
\;\longmapsto\;
\left[
\begin{array}{c}
\hat{v}_0 \\[1pt]
\vdots \\[1pt]
\hat{v}_{n-1} \\[1pt]
\hat{v}_n \\[1pt]
\vdots \\[1pt]
\hat{v}_{511}
\end{array}
\right]
\label{eq:condinput}
\end{equation}
%
from $\mathbb{R}^{512\times 81}$ to $\mathbb{R}^{512\times 40}$, the input
blocks being the interpolant, the window on the given frames and the indicator,
of widths $40$, $40$ and $1$. Here $x_k$ and $X_k$ are the $40$ channels of
frame $k$ in the interpolant and in the window, and $\hat{v}_k$ is the velocity
returned at that row.

The first $n_b$ rows of the output are discarded. The rest are scored against
the target and averaged,
%
\begin{equation}
\ell_b \;=\; \frac{1}{512-n_b} \sum_{k=n_b}^{511}
\bigl\| \hat{v}^{(b)}_k - \bigl(\epsilon^{(b)}_k - X^{(b)}_k\bigr) \bigr\|^2 ,
\label{eq:windowloss}
\end{equation}
%
with the norm the mean squared entry over the forty channels of a row, and the
objective is their mean over the batch,
%
\begin{equation}
L(\theta) \;=\; \frac{1}{32}\sum_{b=1}^{32} \ell_b .
\label{eq:condloss}
\end{equation}
%
Dividing by $512-n_b$ weights every window equally, whatever fraction of it is
generated. On a given row the window already appears in the input, so its target
follows from the interpolation and it is excluded.

The interpolant is noised at every frame and the given frames are never
substituted into it, so training and generation present the same input.

The range of $n_b$ trains one network across every amount of history, and
$n_b=0$ supplies the unconditional case.

\subsection{The optimiser}

The network defines a parameterised function $v_\theta$ with $P$ trainable
parameters,
\[
\theta=(\theta_1,\ldots,\theta_P)\in\mathbb{R}^P.
\]

For an input $(x,t)$, the forward pass produces the prediction $\hat v$ of
\eqref{eq:condinput}, from which the scalar batch loss $L(\theta)$ of
\eqref{eq:condloss} is evaluated. The loss depends on the parameters through the
intermediate computations of the network,
\begin{equation}
(x,t;\theta)
\longrightarrow z_1
\longrightarrow \cdots
\longrightarrow z_m
\longrightarrow \hat v
\longrightarrow L.
\end{equation}

The dependence of the loss on each parameter is described by the gradient
\begin{equation}
\nabla_\theta L
=
\left[
\frac{\partial L}{\partial\theta_1},
\ldots,
\frac{\partial L}{\partial\theta_P}
\right]
\in\mathbb{R}^P.
\end{equation}

Each component $\partial L/\partial\theta_i$ gives the local sensitivity of the
loss to $\theta_i$, with the remaining parameters held fixed. Backpropagation
computes these derivatives efficiently by applying the chain rule backwards
through the computational graph. Along a single path
\[
\theta_i\rightarrow z_1\rightarrow\cdots\rightarrow z_m\rightarrow L,
\]
the corresponding derivative is
\begin{equation}
\frac{\partial L}{\partial\theta_i}
=
\frac{\partial L}{\partial z_m}
\frac{\partial z_m}{\partial z_{m-1}}
\cdots
\frac{\partial z_1}{\partial\theta_i},
\end{equation}
with contributions accumulated over all paths connecting $\theta_i$ to $L$. The
gradient therefore provides a local direction in parameter space along which the
loss changes.

An optimiser converts these gradients into parameter updates. Writing
\begin{equation}
G_{i,t}
=
\frac{\partial L_t}{\partial\theta_i}
\end{equation}
for the gradient of parameter $\theta_i$ at optimisation step $t$, gradient
descent gives
\begin{equation}
\theta_i^{(t+1)}
=
\theta_i^{(t)}
-
\eta G_{i,t},
\end{equation}
where $\eta>0$ is the learning rate. A single learning rate therefore scales all
parameter updates, despite potentially substantial differences in gradient scale
between parameters and over the course of training.

Adam~\cite{adam} instead constructs a parameter-wise adaptive update from the
recent history of each parameter's gradients. It maintains exponentially
weighted estimates of the gradient and squared gradient,
\begin{equation}
M_{i,t}
=
\beta_1 M_{i,t-1}
+
(1-\beta_1)G_{i,t},
\end{equation}
and
\begin{equation}
V_{i,t}
=
\beta_2 V_{i,t-1}
+
(1-\beta_2)G_{i,t}^2.
\end{equation}

Here, $M_{i,t}$ is a smoothed estimate of the recent gradient, while $V_{i,t}$
is a smoothed estimate of its squared magnitude. Consequently,
$\sqrt{V_{i,t}}$ provides an estimate of the recent root-mean-square gradient
magnitude. Dividing the first moment by this quantity normalises the update by
the characteristic gradient scale of that parameter.

Both moments are initialised at zero, $M_{i,0}=V_{i,0}=0$, which biases their
early values towards zero. Adam removes this initialisation bias using
\begin{equation}
\hat M_{i,t}
=
\frac{M_{i,t}}{1-\beta_1^t},
\qquad
\hat V_{i,t}
=
\frac{V_{i,t}}{1-\beta_2^t}.
\end{equation}
The factors $1-\beta_1^t$ and $1-\beta_2^t$ account for the missing history at
early optimisation steps; as $t$ increases, both corrections approach unity.

The complete Adam update is therefore
\begin{equation}
\theta_i^{(t+1)}
=
\theta_i^{(t)}
-
\eta
\frac{\hat M_{i,t}}
{\sqrt{\hat V_{i,t}} \;+\; \xi},
\end{equation}
where $\xi>0$ ensures numerical stability when the estimated gradient magnitude
is small. The numerator determines the smoothed descent direction, while the
denominator rescales it by the recent gradient magnitude. Adam therefore reduces
the dependence of each update on the absolute scale of its parameter's gradient,
yielding a parameter-wise adaptive effective learning rate.

The optimiser parameters used here are summarised in
Table~\ref{tab:optimiser-settings}.

\begin{table}[H]
\centering\small
\begin{tabular}{lll}
\toprule
Symbol & Meaning & Value \\
\midrule
$\eta$ & learning rate & $5\times10^{-4}$ \\
$\beta_1$ & first-moment decay & $0.9$ \\
$\beta_2$ & second-moment decay & $0.999$ \\
$\xi$ & numerical stability constant & $10^{-8}$ \\
\bottomrule
\end{tabular}
\caption{Optimiser settings.}
\label{tab:optimiser-settings}
\end{table}

Before the Adam moments are formed, the full gradient is clipped to unit norm.
When clipping is active, this preserves the gradient direction while bounding its
magnitude and limiting the effect of anomalously large gradients. Training is
performed for $10^4$ optimisation steps.

Backpropagation thus computes the sensitivity of the loss to each network
parameter, while Adam converts these gradients and their recent first- and
second-moment estimates into parameter updates.

\begin{algorithm}[H]
\caption{Pretraining}
\label{alg:pretrain}
\begin{algorithmic}[1]
\Repeat
  \State $X \sim p$, \; $\epsilon \sim \mathcal{N}(0,I)$, \;
         $t \sim \mathcal{U}[0,1]$, \; $n \sim \mathcal{U}\{0,\ldots,256\}$
  \State $x \gets (1-t)X + t\epsilon$
  \State assemble the input of \eqref{eq:condinput} and read off $\hat{v}$
  \State take an Adam step on $\nabla_\theta \ell$, the error of
         \eqref{eq:windowloss} over rows $n$ to $511$
\Until{converged}
\end{algorithmic}
\end{algorithm}

% ================================================================
\section{Sampling}

At inference, sampling begins from a $512\times40$ tensor of independent
standard Gaussian noise,
\[
x_1=\varepsilon,
\qquad
\varepsilon_{ij}\overset{\mathrm{i.i.d.}}{\sim}\mathcal N(0,1).
\]

The network $v_\theta(x_t,t;c)$ parameterises a time-dependent velocity field
over the state space. This field describes how the current state $x_t$ should
change at flow time $t$ and defines the ODE
\[
\frac{dx_t}{dt}=v_\theta(x_t,t;c),
\]
where $c$ denotes the conditioning information of \eqref{eq:condinput}, the
given frames together with their indicator. Sampling solves this ODE from $t=1$
to $t=0$, transporting the initial Gaussian noise towards the learned data
distribution.

The flow coordinate $t$ is artificial and is distinct from market time. Market
time is represented by the 512 rows of $x_t$, separated by $1\,\mathrm{ms}$.
Each integration step therefore updates the complete $512\times40$ window
jointly rather than generating market frames sequentially. The temporal
evolution of the order book is learned as part of the joint structure of the
complete window.

For a context length $n$, the observed prefix $X_{<n}$ is supplied through the
conditioning channels of \eqref{eq:condinput}, while the remaining $512-n$
frames are generated. The resulting continuation approximates a sample from
\[
X_{\geq n}
\sim
p\!\left(X_{\geq n}\mid X_{<n}\right).
\]

Using Euler integration with step size $h$, the state is updated according to
\[
x_{t-h}
=
x_t-h\,v_\theta(x_t,t;c),
\]
with the velocity field re-evaluated after each step. At $t=0$, the generated
rows $X_{\geq n}$ are retained and transformed back to the original order-book
representation. Independent noise draws produce different plausible
continuations from the same observed context.

When $n=0$, no market context is supplied and the model generates an
unconditional $512\,\mathrm{ms}$ window. For $n>0$, the generated horizon is
$512-n$ milliseconds. The model therefore conditions only on the observed prefix
within the current window, while earlier market history is not explicitly
available during generation.


\begin{thebibliography}{99}\small

\bibitem{bookdit} \textbf{BookDiT.}
\newblock Chart representation and WanDiT backbone for limit order book
diffusion.
\newblock ICAIF 2026 submission, cited by \texttt{src/state}.

\bibitem{adam} \textbf{Adam.} D. P. Kingma, J. Ba.
\newblock Adam, A Method for Stochastic Optimization.
\newblock ICLR 2015.
\href{https://arxiv.org/abs/1412.6980}{arXiv 1412.6980}.

\bibitem{rectified} \textbf{Rectified flow.} X. Liu, C. Gong, Q. Liu.
\newblock Flow Straight and Fast, Learning to Generate and Transfer Data with
Rectified Flow.
\newblock ICLR 2023, UT Austin.
\href{https://arxiv.org/abs/2209.03003}{arXiv 2209.03003}.

\end{thebibliography}

\end{document}
